\chapter{The Exponential Time Hypothesis}
\label{ch:eth}

$\Pc \ne \NP$ says that satisfiability has no polynomial time algorithm. It does not say how much time satisfiability needs, and for thirty years after Cook the question of the exponent was left to the algorithm designers, who brought the time for $3$ satisfiability on $n$ variables down from $2^n$ to $(4/3)^n$ without any theory of where the descent would stop. The exponential time hypothesis\index{exponential time hypothesis} of Impagliazzo and Paturi \cite{IP2001}, formulated in 1999 and sharpened in 2001 with Zane \cite{IPZ2001}, supplied the theory. It asserts that the descent stops at an exponential: $3$ satisfiability needs time $2^{\delta n}$ for some $\delta > 0$. The hypothesis turned out to be the right strength to transfer. Through the sparsification lemma and a web of fine grained reductions it gives tight lower bounds for hundreds of problems, parameterised and counting, and under its strong form polynomial, and it has become the working assumption of a field that $\Pc \ne \NP$ alone could not organise. This chapter states the hypothesis and its variants, proves the consequence for the book's two halves, and asks what Monte Carlo methods can and cannot do against an exponential that is assumed rather than feared.

\section{The hypothesis}

Let $s_k$ be the infimum of the $\delta$ such that $k$ satisfiability on $n$ variables can be decided in time $2^{\delta n}$ by a deterministic algorithm.

\begin{definition}
The exponential time hypothesis, ETH, is the statement $s_3 > 0$: there is a $\delta > 0$ such that $3$ satisfiability cannot be decided in time $2^{\delta n}$. The strong exponential time hypothesis, SETH, is the statement $\lim_{k \to \infty} s_k = 1$: for every $\epsilon > 0$ there is a $k$ such that $k$ satisfiability cannot be decided in time $2^{(1 - \epsilon) n}$.
\end{definition}

ETH implies $\Pc \ne \NP$ and is stronger; SETH implies ETH and is stronger still, and it is believed less firmly, since it asserts that the trivial algorithm is essentially optimal for large $k$ and several of its consequences are tight bounds for polynomial problems that would be surprising if false but are not known to be true. Impagliazzo and Paturi \cite{IP2001} proved that $s_k \le (1 - d/k) s_\infty$ for a constant $d > 0$, so that under ETH the nondecreasing sequence $s_k$ increases infinitely often.

The tool that makes ETH transfer is the sparsification lemma of Impagliazzo, Paturi and Zane~\cite{IPZ2001}.

\begin{theorem}[Sparsification]
For every $\epsilon > 0$ and $k$ there is a constant $c = c(k, \epsilon)$ such that every $k$ CNF formula on $n$ variables can be written, in time $2^{\epsilon n} \mathrm{poly}(n)$, as a disjunction of at most $2^{\epsilon n}$ $k$ CNF formulas each with at most $c n$ clauses.
\end{theorem}

The lemma lets the hypothesis be stated in terms of $n + m$, variables plus clauses, rather than $n$ alone: under ETH there is no $2^{o(n + m)}$ algorithm for $3$ satisfiability. Since the standard reductions from satisfiability produce instances of size linear in $n + m$, every $\NP$ complete problem whose reduction is linear inherits the bound that no $2^{o(\text{size})}$ algorithm exists, and every problem whose reduction is polynomial of degree $d$ inherits the bound that no $2^{o(\text{size}^{1/d})}$ algorithm exists. This is the discovery: that the exponent survives reductions in a controlled way, so that a single hypothesis about one problem fixes the exponential complexity of all of them up to the degree of the reduction.

\section{What ETH implies}

A sample of the consequences, taken from the survey of Lokshtanov, Marx and Saurabh \cite{LMS} and the book of Cygan and coauthors \cite{Cygan}, shows the reach. Under ETH there is no $2^{o(n)}$ algorithm for independent set, vertex cover, dominating set, Hamiltonian cycle or $3$ colouring on $n$ vertex graphs; no $2^{o(\sqrt n)}$ algorithm for the planar versions, which is tight since planar separators give $2^{O(\sqrt n)}$; no $f(k) n^{o(k)}$ algorithm for clique with parameter $k$, which is tight against the trivial $n^k$; no $2^{o(k)} \mathrm{poly}(n)$ algorithm for vertex cover parameterised by the solution size $k$. Under SETH, Williams \cite{Williams2005} and others showed that orthogonal vectors has no $n^{2 - \epsilon}$ algorithm, hence neither does edit distance, and that dominating set has no $n^{k - \epsilon}$ algorithm; these are the fine grained lower bounds, and they say that polynomial problems too have exponents fixed by the hypothesis.

For counting there is a counting version. Dell, Husfeldt, Marx, Taslaman and Wahl\'en \cite{DHMTW} formulated $\sharp$ETH, that $\sharp 3$SAT has no $2^{o(n)}$ algorithm, and proved under it that the permanent of an $n \times n$ matrix cannot be computed in time $2^{o(n)}$, so that Ryser's $2^n$ algorithm is optimal up to the exponent, and that the Tutte polynomial has exponential lower bounds at almost every point. The counting hypothesis is weaker than ETH, since it follows from ETH because counting decides; its use is that it reaches problems like the permanent whose decision version is easy, where no decision reduction from satisfiability is available.

For the queens, the reduction of Gent, Jefferson and Nightingale of Chapter~\ref{ch:npcomplete} produces a board of size polynomial in the formula, and we write $N = \Theta((n + m)^d)$ for its degree. Under ETH the completion problem has no algorithm of time $2^{o(N^{1/d})}$, and under $\sharp$ETH the same holds for counting completions; a linear reduction, if one exists, would give $2^{\Omega(N)}$, which is far below the $N! \approx 2^{N \log N}$ of exhaustive search and far above polynomial. The unconditioned count $Q(N)$ is untouched by any of this, since it is a function of one integer and no hypothesis about formulas is known to constrain a single sequence; the realistic prize of Chapter~\ref{ch:program} remains available under ETH.

\section{ETH and the two halves of the program}

ETH is a hypothesis about deterministic time, and the book's thesis is about randomized time. The two meet through the circuit hypothesis of Impagliazzo and Wigderson \cite{IW} and through the randomized form of ETH.

\begin{theorem}[ETH and the program]
\label{thm:ethprogram}
\begin{enumerate}
\item If the class $\mathrm E$ of problems decidable in time $2^{O(n)}$ contains a problem with no circuits of size $2^{\epsilon n}$ for some $\epsilon > 0$, then $\Pc = \BPP$, hence $\Pc = \RP$: the derandomization half of the program holds.
\item Assume ETH and the hypothesis of (1). Then $\Pc = \RP = \BPP \ne \NP$, and consequently the queens level sets on the boards of the Gent Jefferson Nightingale reduction admit no polynomial sampler, and no fully polynomial randomized approximation scheme for completion counts exists. The Monte Carlo half fails on hard boards for samplers and approximation schemes; we do not prove the corresponding statement for particular realisations of the vertical walk.
\item Assume the randomized exponential time hypothesis, that there is $\delta > 0$ such that $3$ satisfiability has no randomized algorithm of time $2^{\delta n}$. Then on the reduction boards of a $3$ CNF formula with $n$ variables, every randomized algorithm that samples completions within a polynomial factor of uniform and every FPRAS for the completion count has running time $2^{\Omega(n)}$ for infinitely many $n$.
\end{enumerate}
\end{theorem}

\begin{proof}
(1) is the theorem of Impagliazzo and Wigderson of Chapter~\ref{ch:pseudo}, and $\RP \subseteq \BPP$. (2): ETH gives $\Pc \ne \NP$, and with $\Pc = \RP$ this is $\RP \ne \NP$; the main equivalence of Chapter~\ref{ch:theorems} then denies the sampler and the scheme. (3): a polynomial factor sampler on the reduction board of a formula yields, by the main equivalence and the bootstrapping of the polynomial factor of Sinclair and Jerrum \cite{SinclairJerrum}, a randomized algorithm for the formula's satisfiability in time polynomial in $N$, hence polynomial in $n$; a randomized polynomial algorithm for $3$ satisfiability contradicts the randomized hypothesis. Moreover the contradiction is quantitative: a sampler of time $T(N)$ gives a satisfiability algorithm of time $T(N) \mathrm{poly}(N)$, so $T(N) \ge 2^{\delta n} / \mathrm{poly}(N)$, which is $2^{\Omega(n)}$. The same holds for an FPRAS for the completion count, since it separates count $0$ from count at least $1$ and so decides the formula's satisfiability in randomized time polynomial in its running time and $N$.
\end{proof}

The theorem says that the world the book expects, $\Pc = \RP \ne \NP$, is a consequence of the exponential time hypothesis together with the circuit hypothesis of Impagliazzo and Wigderson, both of which are believed by most of the field, and that under the randomized form the Monte Carlo half fails on hard boards not merely in the sense of not being polynomial but in the sense of being exponential in the formula size for every sampler and approximation scheme (whether the vertical walk in particular is covered we do not prove). That is the honest status of the program's Monte Carlo half under the randomized form of the hypothesis: it is a polynomial method on the boards that occur and an exponential one on the boards that are built.

\section{Where the exponential must sit}

The theorems of Part IV locate the exponential precisely, and ETH says it is there to be located. The vertical walk's cost is a product of three factors: the number of rungs, which is polynomial by the integer ladder; the quadrature cost per rung, which for absolute error is $O(1 / \epsilon)$ nodes with an absolute constant by the staircase theorems of Chapter~\ref{ch:staircase}; and the cost of a level set sample, which is the mixing time of the chain or, with rung skipping, the mixing time at the first connected level times the inverse ratio of the ground state to that level. Under the randomized hypothesis, for a realisation of the walk that is a sampler or an FPRAS in the sense of Theorem~\ref{thm:ethprogram}(3), the whole product is $2^{\Omega(n)}$ for infinitely many $n$, on the reduction boards of hard formulas (the board of an easy formula is easy), and, if the first two factors are polynomial (the second only given a lower bound on the rung integral, which we do not prove, since the conversion from absolute to relative error carries a factor $1/Z$, Chapter~\ref{ch:staircase}), the third is exponential in aggregate. The level by level form below is a heuristic that we do not prove, since an algorithm that uses level $m$ also needs a starting state in $\{S \le m\}$ and the contradiction bounds the whole walk rather than each level separately: on such boards, for the levels $m$ the walk uses,
\[
\tau_m \cdot \frac{ |\{ S \le m \}| }{ Q(N; X_0) } \ge 2^{\Omega(n)} ,
\]
where $\tau_m$ is the mixing time of any local chain on the conditioned level set $\{ S \le m \}$. Heuristically, either the chain is torpid at the levels whose ratio to the ground state is subexponential, or the levels at which the chain mixes have exponentially more placements than completions. This is a weak conditional relative of the shattering depth conjecture of Chapter~\ref{ch:open}, with the depth measured not in attacks but in the logarithm of the ratio, and it is the one statement about the queens landscape on hard boards that the book can make without computing anything, though we have not proved it.

The constant in the exponent is a different matter. The randomized hypothesis gives $2^{\delta n}$ with $\delta$ unknown and small; the best algorithms for $3$ satisfiability, which are Monte Carlo, give $\delta \le \log_2 1.308 \approx 0.387$. These algorithms are the subject of the next section, because they are level set chains.

\section{Monte Carlo at exponential scale}

The best known algorithms for $3$ satisfiability are randomized, and the first of them is a local search, a level set chain of exactly the kind this book studies. Sch\"oning's algorithm \cite{Schoning} of 1999 starts from a uniformly random assignment, and for $3n$ steps picks an unsatisfied clause, picks one of its $k$ literals uniformly, and flips it; it restarts if no satisfying assignment is found. The energy is the number of unsatisfied clauses, the moves are single flips chosen among those that touch a violated constraint, and the analysis is a one dimensional random walk on the Hamming distance to a fixed satisfying assignment: each step decreases the distance with probability at least $1/k$, since at least one literal of a violated clause is set the wrong way, and a walk of $O(n)$ steps reaches distance zero with probability at least $(1/(k-1))^{d}$ from distance $d$ up to polynomial factors, which averaged over the random start gives success probability $(k / (2(k-1)))^n$ per restart and expected time $(2(k-1)/k)^n$, equal to $(4/3)^n$ for $k = 3$. Paturi, Pudl\'ak, Saks and Zane's algorithm \cite{PPSZ}, which chooses variables in a random order and sets each by unit resolution when it can and at random when it cannot, achieves $1.308^n$ after Hertli's analysis \cite{Hertli}, and it is a Monte Carlo method too. Under ETH these constants cannot be driven to one; under SETH the bases for $k$ satisfiability must tend to two as $k$ grows. The Monte Carlo methods of the satisfiability community are the exponential time regime of the book's level set chains, with the exponent as the object of study rather than the polynomial.

The same analysis can be written for the queens, and it is instructive because it fails.

\begin{proposition}[A Sch\"oning walk for the queens]
\label{prop:schoening_queens}
Consider the chain on placements of a board with $f \ge 2$ free rows (permutations of the free columns) which, from a placement with at least one attacking pair, picks an attacking pair uniformly, picks one of its two queens uniformly, the free one if the other is fixed, and swaps that queen's row with a uniformly random other free row. Started from a uniform random placement of the free rows and run for $f$ steps with restarts, the chain finds a completion, if one exists, in expected time at most $(2 f)^{f} \cdot \mathrm{poly}(f)$.
\end{proposition}

\begin{proof}
Fix a completion $x^*$ and let $d$ be the number of free rows on which the current placement differs from $x^*$. Of the two queens of an attacking pair at least one is on a row where the placement differs from $x^*$, since two queens both placed as in $x^*$ do not attack; with probability at least $1/2$ the chain picks such a queen, and with probability $1/(f - 1)$ it swaps its row with the row that currently holds the column $x^*$ assigns to it, after which both rows agree with $x^*$ or the distance has decreased by at least one. So each step decreases $d$ with probability at least $1 / (2 (f - 1))$, and the walk reaches $d = 0$ within $d$ steps with probability at least $(2 f)^{-d} \ge (2 f)^{-f}$. Restarting $(2 f)^f$ times succeeds with constant probability.
\end{proof}

The bound $(2f)^f$ is worse than exhaustive search over the $f!$ placements of the free rows. The reason is the one that separates the queens from satisfiability: a violated constraint in a $k$ CNF formula points at $k$ variables, at least one of which is wrong, while an attacking pair points at two queens, at least one of which is wrong, but says nothing about where the wrong queen should go, and the choice among $f - 1$ destinations is what costs the factor $f$ per step. Sch\"oning's method gains over exhaustive search because the constraint width is bounded; the queens have constraints of width two and a destination alphabet of size $N$, and the trivial bound gives nothing. What is not trivial is the level set structure: the chain of Chapter~\ref{ch:experiments2}, which accepts any swap keeping at most one attack rather than only swaps that repair the chosen pair, has polynomial autocorrelation time on every board tried, and the gap between $(2f)^f$ and polynomial is the gap between an analysis that tracks one Hamming distance and the geometry of the landscape that the book has been measuring. The honest statement is that the analysis above gives no exponent below $f \log f$, that the measured behaviour is polynomial on random boards, and that the randomized form of ETH, through Theorem~\ref{thm:ethprogram}(3), says that a sampler or approximation scheme built on the measured behaviour must give way to $2^{\Omega(n)}$ on the reduction boards of hard formulas, for infinitely many $n$, without saying at which level set, which is what the shattering depth would tell.

\section{How the Monte Carlo methods tackle ETH hard problems}

Three ways, each with its limit stated by the hypothesis.

Average case. ETH and its relatives are worst case statements, and the walk on random and planted boards is an average case algorithm. Nothing in the hypothesis prevents a problem that is $2^{\Omega(n)}$ in the worst case from being polynomial on every instance a user generates, and the shattering results of Chapter~\ref{ch:open} are the beginning of an account of where the worst cases are. Levin's theory of average case completeness, in Chapter~\ref{ch:ogp}, says that some distributions are as hard as the worst case; it does not say that the uniform distribution on boards with $k$ queens is one of them, and the evidence of this book is that for $k$ below the threshold it is not. The limit is that the hard instances exist, by ETH, and the walk will meet them wherever a formula is encoded.

Moderately exponential. Under the randomized hypothesis the cost on hard boards of any sampler or approximation scheme, and heuristically of the walk, is $2^{\Omega(n)}$ for infinitely many $n$ by Theorem~\ref{thm:ethprogram}(3), but the constant is the question, and Sch\"oning's and PPSZ's algorithms show that Monte Carlo methods set the record for the constant on satisfiability. The vertical walk's analogue is a level set chain on a reduction board whose mixing time is $2^{cN}$ with $c$ small, and the product estimator across rungs then gives a counting algorithm of the same exponential cost, which for counting completions would be competitive with the $2^{n}$ of exhaustive enumeration of the formula if $c N$ is below $n$. Whether any chain achieves this is open; the exact spectral computations of Chapter~\ref{ch:open} are the instrument that would measure $c$ on small reduction boards.

Approximation and derandomization. ETH is a hypothesis about exact decision; the walk is an approximate counter. The main equivalence says that the two meet on self reducible problems, where approximate counting is as hard as decision, so approximation buys nothing against ETH for the queens. But the circuit hypothesis of Impagliazzo and Wigderson, by Theorem~\ref{thm:ethprogram}(1), gives $\Pc = \RP$, and that is the one place where a hardness hypothesis helps the program rather than limiting it: the derandomization half of the thesis follows from it. The reader who believes ETH and the circuit hypothesis believes that every polynomial walk the program might find can be run without a coin, and that no polynomial walk exists on reduction boards. That is the book's expected world, derived from these two hypotheses.

\section{Summary}

The exponential time hypothesis gives the book's program its scale. Under ETH the completion problem is exponential in a power of $N$ and the unconditioned count is unconstrained; under the circuit hypothesis of Impagliazzo and Wigderson $\Pc = \RP$ and the derandomization half follows; under the randomized form every sampler and approximation scheme, and heuristically every version of the vertical walk, is $2^{\Omega(n)}$ on the boards that encode a hard formula, for infinitely many $n$, and the exponential sits entirely in the sampling, as, heuristically, a torpid chain or an exponentially small ratio. Monte Carlo methods tackle such problems by working on average, where they are polynomial on every board the book has tried; by setting the constant in the exponent, where they hold the records for satisfiability and where the queens analogue is open; and by not needing to be derandomized at all, if the hypotheses that make them fail in the worst case also make their coins unnecessary.
